%% LyX 2.4.4 created this file.  For more info, see https://www.lyx.org/.
%% Do not edit unless you really know what you are doing.
\documentclass[english]{article}
\usepackage[T1]{fontenc}
\usepackage[latin9]{inputenc}
\usepackage{enumitem}
\usepackage{amsmath}
\usepackage{amsthm}
\usepackage{cancel}
\usepackage{geometry}
\geometry{verbose,tmargin=1in,bmargin=1in,lmargin=1in,rmargin=1in}

\makeatletter

%%%%%%%%%%%%%%%%%%%%%%%%%%%%%% LyX specific LaTeX commands.
%% Because html converters don't know tabularnewline
\providecommand{\tabularnewline}{\\}

%%%%%%%%%%%%%%%%%%%%%%%%%%%%%% Textclass specific LaTeX commands.
\numberwithin{equation}{section}
\numberwithin{figure}{section}
\newlength{\lyxlabelwidth}      % auxiliary length 

%%%%%%%%%%%%%%%%%%%%%%%%%%%%%% User specified LaTeX commands.
\usepackage{fancyhdr}
\usepackage{lastpage}
\pagestyle{fancy}
\usepackage{enumitem}
\usepackage{ifthen}
%sets math fonts to be more readable
\everymath=\expandafter{\the\everymath\displaystyle}
%\DeclareMathSizes{10}{10}{10}{10}
\usepackage{multicol}

\usepackage{tikz}
\usetikzlibrary{plotmarks}
\usepackage{pgfplots}
\pgfplotsset{compat=newest}

\usepackage{calc}
\usepackage[nomessages]{fp}

\usepackage{advdate}    % Advancing/saving dates
\usepackage{datetime}   % Dates formatting
\usepackage{datenumber} % Counters for dates
% Added by lyx2lyx
\setlength{\parskip}{\smallskipamount}
\setlength{\parindent}{0pt}

\makeatother

\usepackage{babel}
\begin{document}
\renewcommand{\author}{Dr.\,James Ashe}

\newcommand{\classNum}{107}

\newcommand{\classSec}{7 and 10}

\newcommand{\nameType}{\hspace{2cm}Name:}

\newcommand{\examType}{Quiz 3: 1.5}

\renewcommand{\date}{\formatdate{15}{9}{2014}}

%%First page's header%%%%%%%%%%%%%%%%%%%%%%
\fancypagestyle{firststyle} %%header settings for first pagr
{
	\fancyhead[L]{MTH \classNum{}(\classSec{}) \author{}\\
	\textbf{\examType{}} \date{}}
	\fancyhead[C]{\textbf{\nameType}}
	\setlength{\headsep}{14pt}
}
%%%%%%%%%%%%%%%%%%%%%%%%%%%%%%%%%%%%%%%%%%

%%Next page's header%%%%%%%%%%%%%%%%%%%%%%%
%\setlist{nolistsep}
\lhead{MTH \classNum{}(\classSec{})} 
\chead{\textbf{\nameType}} 
\cfoot{\thepage{}~of~\pageref{LastPage}} 
\setlength{\headsep}{5pt} 
%%%%%%%%%%%%%%%%%%%%%%%%%%%%%%%%%%%%%%%%%%%

%%Various settings; see the preamble for other settings%%%%%
%%\setenumerate[0]{leftmargin=12pt,itemindent=0pt,labelwidth=0pt}
\renewcommand{\labelenumi}{\textbf{\arabic{enumi}.}}
\renewcommand{\labelenumii}{\textbf{\alph{enumii}.}}
\thispagestyle{firststyle}
%%%%%%%%%%%%%%%%%%%%%%%%%%%%%%%%%%%%%%%%%%
\newcommand{\xmax}{10}
\newcommand{\xmin}{-10}
\newcommand{\xstep}{0} %must be xmin+n
\newcommand{\ymax}{10}
\newcommand{\ymin}{-10}
\newcommand{\ystep}{0} %must be ymin+n

\newcommand{\xspace}{\vspace{.75cm}}

\emph{You must show work to receive credit. Unless indicated otherwise,
each problem is worth equal value.}
\begin{enumerate}[resume]
\item Construct a Venn diagram to determine the validity of the given argument.

\begin{enumerate}
\item \vspace{-25pt}%
\begin{tabular}{l}
\tabularnewline
\tabularnewline
1. All homeless people are unemployed.\tabularnewline
2. Bill Gates is not unemployed.\tabularnewline
\hline 
Therefore, Bill Gates is not homeless\tabularnewline
\end{tabular}\xspace\xspace
\item \vspace{-25pt}%
\begin{tabular}{l}
\tabularnewline
\tabularnewline
1. No vegetarian owns a gun.\tabularnewline
2. All policemen own guns.\tabularnewline
\hline 
Therefore, no policeman is a vegetarian.\tabularnewline
\end{tabular}\xspace\xspace
\item \vspace{-25pt}%
\begin{tabular}{l}
\tabularnewline
\tabularnewline
1. No professor is a millionaire.\tabularnewline
2. No millionaire is illiterate.\tabularnewline
\hline 
Therefore, no professor is illiterate.\tabularnewline
\end{tabular}\xspace\xspace\vfill
\end{enumerate}
\end{enumerate}
\begin{enumerate}
\item Which of the following are statements? Why or why not?

\begin{enumerate}
\item $3+7=9$\xspace
\item Solve the equation $3+7=x$.\xspace
\item Is $\sqrt{2}$ an irrational number?\xspace\vfill
\end{enumerate}
\item Write a sentence that negates the statement.

\begin{enumerate}
\item Some cats have three legs.\xspace
\item No cat has three legs. \xspace
\item The cat is not blue.\xspace\vfill\newpage
\end{enumerate}
\item Translate the sentence into symbolic form.

\begin{enumerate}
\item It is a square and a rectangle. \xspace
\item If John is male then he is human.\xspace
\item All squares are rectangles.\xspace
\item All whole numbers are even or odd.\xspace
\item No square is a triangle.\xspace
\item Being an orthodontist is sufficient for being a dentist.\xspace
\item Not being a monkey is necessary for being an ape.\xspace\vfill
\end{enumerate}
\item Using the symbolic representations

\begin{itemize}
\item []$p$: The car costs $\$70,000$.
\item []$q$: The car goes $140$ mph.
\item []$r$: the car is red.
\end{itemize}
express the following compound statements in symbolic form.
\begin{enumerate}
\item All red cars go $140$ mph.\xspace
\item The car is red, goes $140$ mph, and does not cost $\$70,000$.\xspace
\item If the car does not cost $\$70,000$, it does not cost $\$70,000$.\xspace
\item The car is red, and it does not cost $\$70,000$ or go $140$ mph.\xspace
\item Being able to $140$ mph is sufficient for a car to cost $\$70,000$
or be red.\xspace
\item Not being red is necessary for a car to cost $\$70,000$ and not go
$140$ mph.\xspace\vfill
\end{enumerate}
\end{enumerate}

\end{document}
